\documentclass[10pt]{article}
\usepackage[margin=2.0cm]{geometry}
\usepackage{graphicx}\usepackage{booktabs}\usepackage{amsmath}\usepackage{amssymb}
\usepackage{xcolor}\usepackage[hidelinks]{hyperref}\usepackage{enumitem}\usepackage{fontspec}
\setlength{\parskip}{4pt}\setlength{\parindent}{0pt}\setlist{itemsep=1pt,topsep=2pt}
\graphicspath{{../../../output/08_two_color/cambaliza2000/pulses/delay_sweep/}}
\title{\vspace{-1.2cm}\textbf{Review 07. Two-color excitation with optical pulses and a delay:\\ derivations and checks against the four papers}\\
\large \texttt{mcdo\_2ce} \quad(\texttt{mcdo/twobeam.py}, \texttt{mcdo/pulsed2ce.py},
\texttt{scripts/cambaliza2000\_pulses\_*.py})}
\author{}\date{2026-11-10}
\begin{document}\maketitle\vspace{-1.0cm}

\section*{0. Outcome of the review}
I read the four papers in full, compared each equation with the code, derived each step of the pulse and delay
results again, and tested each step numerically. The papers are Cambaliza \& Saloma, Opt.\ Commun.\ 184, 25--35 (2000);
Lim \& Saloma, Opt.\ Commun.\ 207, 121--130 (2002); Lim \& Saloma, Appl.\ Opt.\ 42(17), 3398--3406 (2003); and
Romallosa, Bantang \& Saloma, Phys.\ Rev.\ A 68, 033812 (2003).
\begin{enumerate}
\item The scalar field, the 2CE and 2PE signals, the annular pupil, the misalignment, the aberration function and the
Linfoot criteria in the code are those of the papers. Section 1 lists each equation and its place in the code.
\item The vector field had one error. Its longitudinal component carried a factor $-2i$ where Richards \& Wolf and
Romallosa et al.\ (Eq.~4) have $-2$, so $\nabla\cdot\mathbf{E}$ was not zero. I corrected it and added a test.
It changes one stored case, 2PE with crossed lenses and the polarization in the plane of the two beams. It does not
change the delay results or the pulse results, which use the polarization perpendicular to that plane, where
$E_z = 0$ on the points computed.
\item The delay results follow from three steps, each checked: the temporal overlap of two Gaussian pulses,
$\exp(-2\ln 2\,\Delta t^2/\tau^2)$; the point where two facing pulses meet, $z = c\Delta t/2n$; and the phase
that the delay gives to each kind of photon pair. Section 2 has the derivations and Section 3 the checks.
\item Three statements in the delay deck needed a correction, and I made them (Section 4). The fringes against the
delay have three periods, 2.19, 2.50 and 1.17 fs, and the deck had quoted one value, 2.3 fs.
\item Four limits of the model remain and are stated in Section 5.
\end{enumerate}

\section*{1. The papers and the code}
\subsection*{1.1 Equations}
\begin{center}\small\begin{tabular}{p{4.3cm}p{6.6cm}p{5.2cm}}\toprule
Paper & Statement & Code\\\midrule
Cambaliza \& Saloma, Eq.~1; Lim \& Saloma (2002), Eq.~2 & Paraxial field of a uniformly illuminated circular pupil
(Lommel functions), $u = (2\pi/\lambda)\mathrm{NA}^2 z$, $v = (2\pi/\lambda)\mathrm{NA}\,x$ &
\texttt{twobeam.\_paraxial}: $U = e^{ikz}\,2\!\int_0^1 J_0(v\rho)e^{-iu\rho^2/2}\rho\,d\rho$, the integral that the
Lommel functions evaluate\\
Cambaliza \& Saloma, Eq.~2 & $F_{2c} = \kappa I_1 I_2$ & \texttt{twobeam.two\_color}\\
Cambaliza \& Saloma, Eq.~3 & $F_{2p} = \kappa'|U_1 + U_2|^4$ & \texttt{twobeam.two\_photon}\\
Cambaliza \& Saloma, Eqs.~4--6 & Annular pupil, $U_a = U_R - U_{\varepsilon R}$ & the pupil integral from
$\rho = \varepsilon$ to 1, which is the same difference\\
Lim \& Saloma (2002), Sec.~2 & Focus of beam 2 at $O'(x_c, z_c)$ & arguments \texttt{xc}, \texttt{zc} of
\texttt{beam\_field}\\
Lim \& Saloma (2002), Sec.~3.3; (2003), Sec.~3B & $F = 1 - \langle(t-a)^2\rangle/\langle t^2\rangle$,
$C = \langle a^2\rangle/\langle t^2\rangle$, $Q = \langle|a||t|\rangle/\langle t^2\rangle$, $F = 2Q - C$ &
\texttt{mcdo.linfoot}; the identity is a test\\
Lim \& Saloma (2003), Eq.~1 & Aberration function $\Phi = (\phi/100)\lambda\rho^4$, so the pupil phase is
$k\Phi = 2\pi(\phi/100)\rho^4$ & argument \texttt{sa} of \texttt{beam\_field}\\
Romallosa et al., Eqs.~5--7 & $I_0, I_1, I_2$ with $\cos^{1/2}\vartheta$ and the factors $\sin\vartheta(1+\cos\vartheta)$,
$\sin^2\vartheta$, $\sin\vartheta(1-\cos\vartheta)$ & \texttt{twobeam.\_rw}, the same factors\\
Romallosa et al., Eqs.~2--4 & Components of the field & see 1.2, item 2\\
Romallosa et al., Eqs.~8--9 & $|E(t)|^2 \propto \exp[-4\ln 2\,(t/\tau)^2]$, with $\tau$ the FWHM of the intensity, and
$E(\omega) \propto \exp[-(\omega-\omega_c)^2/4a]$, $a = 2\ln 2/\tau^2$ & \texttt{pulsed2ce.spectrum}:
$a(\omega) = \exp[-(\omega-\omega_0)^2\tau^2/(8\ln 2)]$, the same function\\
\bottomrule\end{tabular}\end{center}
The numbers that Romallosa et al.\ quote follow from these equations. The time--bandwidth product is
$\Delta\nu\,\tau = 4\ln 2/2\pi = 0.441$. A 1-fs pulse at 750 nm has $\Delta\nu = 4.4\times10^{14}$ Hz and spans 483 nm to
1.67 $\mu$m. A 1.2-fs pulse has $c\tau = 359.7$ nm, below $\lambda_c/2 = 375$ nm.

\subsection*{1.2 Points where the current work departs from the papers, or reads them in a particular way}
\begin{enumerate}
\item Cambaliza \& Saloma and Lim \& Saloma (2002) rotate the optical coordinates,
$u_r = u\cos\theta - v\sin\theta$, $v_r = v\cos\theta + u\sin\theta$. Since $u$ scales with $\mathrm{NA}^2$ and $v$ with
NA, this is a rotation of real space at $\theta = 0$ and $\pi$ but not at $\pi/2$. The rotation of real space
reproduces the fits of Lim \& Saloma (2002) ($-0.063\,x_c^2$ against $-0.06\,x_c^2$), and the rotation of the optical
coordinates does not ($-0.21\,x_c^2$), so the replications and the current work rotate real space. The code keeps the
other form as \texttt{rotation='optical'}.
\item Romallosa et al.\ print $E_x = -j(I_0 + 2I_2\cos\phi)$ and $E_y = -jI_2\sin\phi$ (Eqs.~2 and 3). Richards \& Wolf
(1959) have $E_x = -i(I_0 + I_2\cos 2\phi)$ and $E_y = -iI_2\sin 2\phi$, and the code uses these. With them and
$E_z = -2I_1\cos\phi$ (Eq.~4), $\nabla\cdot\mathbf{E} = 0$ to $2\times10^{-8}\,k|E|$ (Section 3).
\item Romallosa et al.\ (Eq.~10) add the intensities of the frequencies,
$|E(P)|^2 = \int|E(\omega)|^2|E(P,\omega)|^2d\omega$. This is the intensity integrated over time, since
$\int|E(P,t)|^2dt = 2\pi\int|\tilde E(P,\omega)|^2d\omega$ (Parseval). For two-photon signals they square this
quantity. The current work keeps the field in time and integrates the two-photon amplitude (2.3), since
$\int I^2dt$ is not $(\int I\,dt)^2$ and since 2CE depends on when the two pulses arrive.
\item The field of each frequency carries the factor $\omega/\omega_0$ of the diffraction integral (the $1/\lambda$ of
Cambaliza \& Saloma's Eq.~1). Romallosa et al.\ do not include it. It changes the 5-fs results by about 2\% (2.6) and
the results at 20 fs and above by less than 0.2\%.
\item With $\lambda_1 = 750$ nm and $\lambda_2 = 656$ nm, $1/\lambda_e = 1/\lambda_1 + 1/\lambda_2$ gives
$\lambda_e = 349.93$ nm. The papers and the code take 350 nm. The level is then $1.1\times10^{-3}$ rad/fs from
$\omega_1 + \omega_2$, which multiplies the 350-nm pairs by a constant that is the same at all delays and positions
(0.9998 in amplitude at 100 fs), so the normalized results do not change.
\item Lim \& Saloma (2003) cite Lim \& Saloma (2002) as pages 111--120. The journal pages are 121--130.
\end{enumerate}

\subsection*{1.3 What the papers say about pulses and timing}
Cambaliza \& Saloma: ``Two-color absorption is possible only in regions of the stained specimen, where $I_1$ and
$I_2$ are simultaneously present in both space and time.'' Lim \& Saloma (2003): ``Achieving optimal temporal
overlap of the 2CE pulses is crucial because spatio-temporal simultaneity of the two confocal beams is a necessary
condition for 2CE''; ``the effect of spectral broadening on the PSF structure is negligible for pulse widths that are
longer than 10 femtoseconds''; and ``fluorescence probes that are used to label the samples are normally selected to
minimize the effects of two-photon excitation from the individual $\lambda_1$ and $\lambda_2$ excitation beams.''
The delay deck quotes the first and third of these sentences as printed. The third is the reason the 2CE signal in
the delay results is the 350-nm pairs.

\section*{2. Derivations}
\subsection*{2.1 The pulse and its spectrum}
The field of a transform-limited Gaussian pulse of carrier $\omega_0$ is
\begin{equation}E(t) = \exp\!\left(-2\ln 2\,\frac{t^2}{\tau^2}\right)e^{-i\omega_0t},\qquad
|E(t)|^2 = \exp\!\left(-4\ln 2\,\frac{t^2}{\tau^2}\right),\end{equation}
so $|E|^2$ falls to one half at $t = \pm\tau/2$ and $\tau$ is its FWHM (Romallosa et al., Eq.~9). With
$\int e^{-at^2}e^{i\Omega t}dt = \sqrt{\pi/a}\,e^{-\Omega^2/4a}$ and $a = 2\ln 2/\tau^2$,
\begin{equation}a(\omega) = \exp\!\left[-\frac{(\omega-\omega_0)^2\tau^2}{8\ln 2}\right]
= \exp\!\left[-\frac{(\omega-\omega_0)^2}{2s^2}\right],\qquad s = \frac{2\sqrt{\ln 2}}{\tau}.\end{equation}

\subsection*{2.2 The pulsed field at the focus}
Each frequency is a continuous wave focused by the lens, with field $\mathbf{E}_\omega(P)$ of Section 1. The pulse at
$P$ is their sum,
\begin{equation}\mathbf{E}(P,t) = \int a(\omega)\,\frac{\omega}{\omega_0}\,\mathbf{E}_\omega(P)\,e^{-i\omega t}\,d\omega .\end{equation}
$\mathbf{E}_\omega$ carries the phase $e^{ikz_b}$ along its beam, so the pulse passes its focus at $t = 0$ and
reaches a point at distance $z_b$ along the beam at $t = nz_b/c$. A delay $\Delta t$ of beam 2 replaces $t$ by
$t - \Delta t$ in its field, carrier included, which multiplies each of its frequencies by $e^{i\omega\Delta t}$.

\subsection*{2.3 The signal: photon pairs that reach the 350-nm level}
For a level at $\omega_e$ that is narrow compared with the spectra, and no level near $\omega_1$ or $\omega_2$, the
two-photon amplitude is the component of $E^2$ at $\omega_e$:
\begin{equation}A(P) = \int E^2(P,t)\,e^{i\omega_et}\,dt = \int\tilde E(P,\omega)\,\tilde E(P,\omega_e-\omega)\,d\omega .\end{equation}
The second form says that each pair of photons with $\omega_a + \omega_b = \omega_e$ contributes, and that the
pairs add with their phases. With $E = E_1 + E_2$,
\begin{equation}A = A^{11} + 2A^{12} + A^{22},\qquad A^{ab} = \int E_a(P,t)\,E_b(P,t)\,e^{i\omega_et}\,dt .\end{equation}
$2A^{12}$ takes one photon from each beam: these are the 350-nm pairs (750 + 656 nm). $A^{11}$ takes both photons
from beam 1, the 375-nm pairs (750 + 750 nm), and $A^{22}$ both from beam 2, the 328-nm pairs (656 + 656 nm). The
signal is $F = |A|^2$; for the vector field the code sums $|A_{ij}|^2$ over the components $i, j$.

\subsection*{2.4 The limit of long pulses is the result of the papers}
For long pulses the spectra are lines. $A^{11}$ and $A^{22}$ oscillate at $\omega_e - 2\omega_1$ and
$\omega_e - 2\omega_2$ and integrate to zero, and $2A^{12} \to 2E_1E_2$, so
\begin{equation}F \to 4|E_1|^2|E_2|^2 \propto I_1I_2\quad\text{(Cambaliza \& Saloma, Eq.~2).}\end{equation}
For 2PE, $\omega_1 = \omega_2 = \omega_e/2$, the three terms are all resonant and
$A \to (E_1 + E_2)^2$, so $F \to |U_1 + U_2|^4$ (their Eq.~3). The stored 100-fs profiles at NA 0.5 match the
continuous-beam ones at the three angles.

\subsection*{2.5 The size of the 375- and 328-nm pairs}
For two Gaussians of the same $s$, $\int a_1(\omega)a_2(\omega_e-\omega)d\omega \propto
\exp[-(\omega_e-\omega_{01}-\omega_{02})^2/4s^2]$. With $4s^2 = 16\ln 2/\tau^2$ and the factor $\omega/\omega_0$ of
each photon evaluated at $\omega_e/2$,
\begin{equation}\frac{|A^{11}|}{|2A^{12}|} = \frac12\left(\frac{\omega_e}{2\omega_1}\right)^2
\exp\!\left[-\frac{(\omega_e-2\omega_1)^2\tau^2}{16\ln 2}\right],\end{equation}
and the same with $\omega_2$ for $A^{22}$. Here $\omega_1 = 2.5115$, $\omega_2 = 2.8714$ and $\omega_e = 5.3819$
rad/fs.
\begin{center}\small\begin{tabular}{lcccc}\toprule
 & \multicolumn{2}{c}{375-nm pairs over 350-nm pairs} & \multicolumn{2}{c}{328-nm pairs over 350-nm pairs}\\
$\tau$ & Eq.~7 & computed at the focus & Eq.~7 & computed at the focus\\\midrule
5 fs & 0.429 & 0.430 & 0.327 & 0.327\\
10 fs & 0.180 & 0.180 & 0.136 & 0.136\\
20 fs & 0.0055 & 0.005 & 0.0040 & 0.004\\\bottomrule\end{tabular}\end{center}
The amplitudes of the pairs from one beam are below 1\% of the 350-nm pairs at 20 fs and reach 18\% at 10 fs. This is
the statement of Lim \& Saloma (2003) that spectral broadening is negligible above 10 fs, with the pulse width at
which it stops being negligible.

\subsection*{2.6 A delay at a fixed point: the temporal overlap}
At a point where both pulses arrive together when $\Delta t = 0$, the envelopes are
$\exp(-2\ln 2\,t^2/\tau^2)$ and $\exp[-2\ln 2\,(t-\Delta t)^2/\tau^2]$. Their product is
\begin{equation}\exp\!\left[-\frac{2\ln 2}{\tau^2}\left(t^2 + (t-\Delta t)^2\right)\right]
= \exp\!\left[-\frac{4\ln 2}{\tau^2}\left(t-\frac{\Delta t}{2}\right)^2\right]
\exp\!\left(-\ln 2\,\frac{\Delta t^2}{\tau^2}\right),\end{equation}
since $t^2 + (t-\Delta t)^2 = 2(t-\Delta t/2)^2 + \Delta t^2/2$. The integral over $t$ of the first factor does
not depend on $\Delta t$, so $|2A^{12}| \propto \exp(-\ln 2\,\Delta t^2/\tau^2)$ and
\begin{equation}\frac{F(\Delta t)}{F(0)} = \exp\!\left(-2\ln 2\,\frac{\Delta t^2}{\tau^2}\right).\end{equation}
The peak is one half at $\Delta t = \tau/\sqrt2 = 0.71\tau$, one quarter at $\Delta t = \tau$ and $0.0039$ at
$2\tau$. The same function is the overlap $\int I_1(t)I_2(t-\Delta t)\,dt$ of the two intensities.
\begin{center}\small\begin{tabular}{lccc}\toprule
Collinear lenses, all NAs & $\Delta t = 0.5\tau$ & $\Delta t = \tau$ & $\Delta t = 2\tau$\\\midrule
Eq.~9 & 0.70711 & 0.25000 & 0.00391\\
stored, 100 fs & 0.70712 & 0.25001 & 0.00391\\
stored, 20 fs & 0.70735 & 0.25034 & 0.00393\\
stored, 5 fs & 0.71095 & 0.25545 & 0.00425\\\bottomrule\end{tabular}\end{center}
The 2\% difference at 5 fs is the factor $\omega/\omega_0$, which makes the spectrum slightly wider than Eq.~2. With
crossed lenses the stored values at 100 fs are 0.2500 to 0.2530.

\subsection*{2.7 Facing lenses: the pulses meet at $c\Delta t/2n$}
Pulse 1 travels along $+z$ and reaches $z$ at $t_1 = nz/c$. Pulse 2 travels along $-z$ with a delay and reaches $z$
at $t_2 = -nz/c + \Delta t$. Eq.~9 holds with $\Delta t$ replaced by the difference of the arrival times,
$t_1 - t_2 = 2nz/c - \Delta t$:
\begin{equation}F_{2c}(z,\Delta t) = I_1(z)\,I_2(z)\,
\exp\!\left[-\frac{2\ln 2}{\tau^2}\left(\frac{2nz}{c}-\Delta t\right)^2\right].\end{equation}
The second factor is largest where the pulses meet, $z = c\Delta t/2n$ (15 $\mu$m in air for 100 fs), and its
FWHM along $z$ is $\sqrt2\,c\tau/2n$ (21.2 $\mu$m at 100 fs, 4.2 $\mu$m at 20 fs). The stored 2CE profile along the
axis is as long as the focus or this stretch, whichever is shorter: its FWHM without delay at 100 fs is 20.7 $\mu$m
at NA 0.1, 15.4 $\mu$m at NA 0.2, 3.3 $\mu$m at NA 0.5 and 0.8 $\mu$m at NA 0.9.

Two limits follow. Where the focus is longer than $c\Delta t/2n$, $I_1I_2$ is nearly constant over the distance the
meeting point moves, so the peak moves with it and is kept. Where the focus is shorter, $I_1I_2$ holds the peak
at the focus and Eq.~9 gives its height.
\begin{center}\small\begin{tabular}{lcccccc}\toprule
Facing lenses, 100 fs, $\Delta t = \tau$ & NA 0.1 & 0.2 & 0.3 & 0.5 & 0.7 & 0.9\\\midrule
peak over the peak without delay & 0.930 & 0.513 & 0.316 & 0.258 & 0.251 & 0.250\\
position of the peak ($\mu$m); $c\tau/2 = 15.0$ & 14.2 & 7.8 & 2.6 & 0.34 & 0.08 & 0.02\\\bottomrule\end{tabular}\end{center}
\begin{center}\includegraphics[width=0.92\textwidth]{delay_facing_vs_na.png}\end{center}
Shorter pulses keep more of the peak at $\Delta t = \tau$, since the meeting point moves by $c\tau/2n$: at NA 0.5 the
peak is 0.258 at 100 fs, 0.409 at 20 fs and 0.868 at 5 fs. Pupils with a longer depth of focus keep more at the same
NA: at NA 0.3 and 100 fs it is 0.316 for the circular pupil, 0.376 for the Gaussian input ($\alpha = 4$) and 0.800
for the annular pupil with $\varepsilon = 0.9$.

\subsection*{2.8 Fringes against the delay}
The delay multiplies each photon of beam 2 by $e^{i\omega\Delta t}$. A 350-nm pair has one such photon, near
$\omega_2$. A 328-nm pair has two, and their frequencies add to $\omega_e$. A 375-nm pair has none. So
\begin{equation}A(\Delta t) = A^{11} + 2A^{12}_0\,g(\Delta t)\,e^{i\omega_2\Delta t} + A^{22}\,e^{i\omega_e\Delta t},\end{equation}
with $g$ the overlap of 2.6. $|A|^2$ then has fringes at the three differences of these phases:
\begin{center}\small\begin{tabular}{llcc}\toprule
Pairs that interfere & Phase & Period & Measured at the focus\\\midrule
350-nm and 375-nm pairs & $\omega_2\Delta t$ & $\lambda_2/c = 2.19$ fs & 2.19 fs\\
350-nm and 328-nm pairs & $(\omega_e-\omega_2)\Delta t$ & $\lambda_1/c = 2.50$ fs & 2.50 fs\\
375-nm and 328-nm pairs & $\omega_e\Delta t$ & $\lambda_e/c = 1.17$ fs & 1.17 fs\\\bottomrule\end{tabular}\end{center}
The first two fall with $g(\Delta t)$. The third keeps its depth at all delays, since neither of its terms needs the
two pulses to overlap. The fringes need pulses with a fixed phase between the two beams; with two lasers that are
not phase-locked they average out and the 350-nm pairs remain (Eq.~9).

In space the same three terms carry the carriers of their beams. With facing lenses these are $e^{2ik_1z}$,
$e^{i(k_1-k_2)z}$ and $e^{-2ik_2z}$, so the 350-nm pairs interfere with either of the others with period
$2\pi/(k_1+k_2) = \lambda_e = 350$ nm along $z$. A delay moves these fringes along $z$. The peak along the axis
follows them, so its fringes against the delay are shallow where the focus is many fringes long and deepen with NA.
With collinear lenses the three terms have the same carrier in space, so the depth is the same at all NAs.
\begin{center}\small\begin{tabular}{lcccc}\toprule
Largest gap, all pairs minus 350-nm pairs & NA 0.1 & 0.3 & 0.5 & 0.7\\\midrule
collinear, 5 fs / 10 fs / 20 fs & \multicolumn{4}{c}{0.917 / 0.710 / 0.036 at all NAs}\\
crossed, along $x$, 5 fs & 0.062 & 0.197 & 0.416 & 0.652\\
crossed, along $x$, 10 fs & 0.064 & 0.191 & 0.413 & 0.599\\
facing, along $z$, 5 fs & 0.014 & 0.013 & 0.033 & 0.070\\
facing, along $z$, 10 fs & 0.008 & 0.015 & 0.035 & 0.076\\\bottomrule\end{tabular}\end{center}
At 20 fs the gap is at most 0.036 in all cases.
\begin{center}\includegraphics[width=0.92\textwidth]{delay_fringe_depth_vs_na.png}\end{center}

\section*{3. Checks}
\begin{center}\small\begin{tabular}{p{7.4cm}p{8.6cm}}\toprule
Check & Result\\\midrule
$\nabla\cdot\mathbf{E}$ of the vector field, NA 0.9, by finite differences & $1.9\times10^{-8}\,k|E|$ with
$E_z = -2I_1\cos\phi$; $0.77\,k|E|$ with the factor $-2i$ that the code had\\
Temporal overlap, Eq.~9, against the stored peaks with collinear lenses & within $1\times10^{-5}$ at 100 fs, 0.006 at
5 fs (table in 2.6)\\
Closed form, Eq.~10, against the stored profiles with facing lenses, NA 0.1 to 0.9, five delays & largest
difference of the normalized profiles 0.004 at 100 fs, 0.019 at 20 fs, 0.039 at 10 fs\\
Ratio of the pairs, Eq.~7, against the amplitudes at the focus & within 0.001 at 5 and 10 fs (table in 2.5)\\
Periods of the fringes against the delay, from the phase of each term at the focus & 2.19, 2.50 and 1.17 fs, as
derived\\
Peak at the edge of the computed line & 187 of 414\,012 stored rows, all at delays of $1.8\tau$ to $2\tau$ where
the peak is below 0.01; none is in a figure\\
Fields from the tables against the original code, before the runs & within $1.6\times10^{-10}$\\
Sixteen samples of the delay runs recomputed with the original code & peaks within $1.8\times10^{-10}$, stored profiles (single precision) within $5.2\times10^{-8}$; all 16 pass\\
Peak over the $x$--$z$ plane against the peak along $x$, crossed lenses & 5 fs, delays 0 to 2.5 fs: the gap is 0.278 over the plane and 0.416 along $x$ at NA 0.5, and 0.478 and 0.652 at NA 0.7\\
Total with all pairs, recomputed as a sum over pairs of frequencies with no time grid (32 cases: both models, the
three angles, 2CE and 2PE, 5 and 10 fs, NA 0.5 to 0.9) & within $1.1\times10^{-4}$ of the peak of the stored profiles\\
350-nm pairs against the pairs summed in frequency and against $\int I_1I_2\,dt$, two other code paths (scalar
model, five cases) & within 0.003 of the peak\\
Test suite & 1074 tests pass, with the corrected $E_z$ and the new test\\\bottomrule\end{tabular}\end{center}
The closed form assumes that each pulse keeps its shape through the focus and that the spot is the same at all
frequencies. The differences of a few percent at 10 fs are the size of these two effects.

\section*{4. Corrections made}
\begin{enumerate}
\item In \texttt{mcdo/twobeam.py} and in its copy in \texttt{scripts/cambaliza2000\_fast\_field.py}, the
longitudinal component is now $E_z = -2I_1\cos\phi$. A new test,
\texttt{test\_vector\_field\_has\_zero\_divergence}, passes with this factor and does not pass with the old one.
The intensity of one beam, the 2CE product $I_1I_2$ and all results with the polarization along $y$ are unchanged.
The 2PE profile with crossed lenses and the polarization in the plane of the beams changes by up to 0.34 of its
peak at NA 0.5 and 0.47 at NA 0.7; it is computed in \texttt{scripts/cambaliza2000\_high\_na\_pupils.py} and I
have not regenerated its outputs.
\item Delay deck, the slides of the fringes: ``one optical period (2.3 fs) apart'' is now ``one optical period
(2.2 to 2.5 fs) apart'', and the notes give the three periods.
\item Delay deck, facing lenses: ``From NA 0.6 the focus is shorter than $c\Delta t/2$'' is now ``At NA 0.5 and
above, the peak stays at the focus and falls with the temporal overlap.'' The focus is shorter than
$c\tau/2 = 15\ \mu$m from about NA 0.2 at 100 fs; the peak is within 0.01 of the temporal overlap from NA 0.5.
\item Delay deck, fringe depth: the notes now say the depth at 20 fs is at most 0.04, and that with crossed lenses
the peak is taken along the $x$ axis through the focus.
\item The reference in the notes is now Appl.\ Opt.\ 42(17), 3398--3406.
\end{enumerate}
\texttt{pulse\_results.pptx} and the main deck are unchanged.

\section*{5. Limits of the model}
\begin{enumerate}
\item The pulses are transform-limited Gaussians and the two beams have a fixed phase. The fringes of 2.8 need
both. The 350-nm pairs alone do not.
\item The level at 350 nm is narrow compared with the spectra and there is no level near 750 or 656 nm. For a broad
absorption band the signal of the 350-nm pairs is $\int I_1(t)I_2(t-\Delta t)\,dt$, which has the same dependence
on the delay (2.6).
\item The sum $\sum_{ij}|A_{ij}|^2$ over the field components is one choice for the response of the molecules. With
the polarization along $y$ on the points computed there is one component, so the choice does not enter the results.
For other polarizations the cross term $2A^{12}_{ij}$ should be made symmetric in $i$ and $j$,
$A^{12}_{ij} + A^{12}_{ji}$, before it is added to $A^{11}$ and $A^{22}$.
\item The peaks are found along the $z$ axis or the $x$ axis through the focus. With facing and collinear lenses
the largest value in the $x$--$z$ plane lies on these lines. With crossed lenses the fringes run along the
diagonal and the peak moves along it, up to 0.25 $\mu$m from the focus at NA 0.5. The fringe depth over the plane is
about two thirds of the depth along $x$ (0.278 against 0.416 at NA 0.5, 0.478 against 0.652 at NA 0.7), and it
deepens with NA in the same way. The crossed-lens values in 2.8 and in the deck are those along $x$.
\end{enumerate}
\end{document}
